% Benchmark section: dataset-benchmark framing, clean channel only.
% Organising principle: every task is a collaborative-perception task, and every
% cohort of baselines is split into LiDAR-only methods and multi-modal methods.
% Communication impairment is deliberately OUT OF SCOPE here (see \S\ref{sec:outlook});
% the only communication axis retained is the message size a method chooses to send,
% which is a property of the method, not of a channel.
% NOTE: \cite keys below are placeholders -- reconcile with the .bib before compiling.

\section{Benchmark Tasks and Baselines}
\label{sec:benchmarks}

The benchmark poses five tasks on the dataset, all of them collaborative: in
each, an \emph{ego} agent must produce an output in its own coordinate frame
that it could not produce as well alone, given messages from the other agents
present in the same scene. Tasks differ in what the output is (geometry, boxes,
pose, identity) and in which modalities may be used to produce it.

Two properties organise the whole section. First, every task is reported
against a \textbf{single-agent reference} computed from the ego's own sensors
with all messages withheld; the collaborative number is meaningful only next to
it. Second, every cohort of baselines is reported in two blocks --- methods
that consume \textbf{LiDAR only}, and \textbf{multi-modal} methods that consume
camera, radar, or CSI in addition to (or instead of) LiDAR. The two blocks are
never merged into one ranking: they have different input budgets, and mixing
them hides whether a gain comes from collaboration or from an extra sensor.

All results in this section are obtained under an \textbf{idealised channel}:
messages are delivered complete, in order, and with zero delay. The dataset
ships logged link measurements, but they are not used here; evaluation under
the logged channel is a separate study and is discussed only in
\S\ref{sec:outlook}.

\subsection{Notation}
\label{sec:bench:notation}

A scene at time $t$ contains agents $\mathcal{A}_t$ with a designated ego
$e \in \mathcal{A}_t$. Agent $i$ observes a LiDAR sweep $P^i_t$, camera images
$\{I^{i,k}_t\}_{k=1}^{K}$, a radar return $R^i_t$ (point cloud and
range--Doppler), and channel state $H^i_t$ (per-subcarrier CSI and RSSI) on its
links, together with a pose $T^i_t \in SE(3)$ in the shared frame. A
collaborative method is a map
\begin{equation}
  \hat{y}^e_t \;=\; F\Bigl(S^e_t,\; \bigl\{\,m^{i \to e}_t \,\bigr\}_{i \in \mathcal{A}_t \setminus \{e\}}\Bigr),
  \qquad
  m^{i \to e}_t \;=\; G\bigl(S^i_t,\, T^i_t,\, T^e_t\bigr),
\end{equation}
where $S^i_t$ is the subset of modalities the method admits, $G$ is its
encoder, and $m^{i \to e}_t$ is the message, whose size $|m^{i \to e}_t|$ in
bits per frame is measured and reported alongside accuracy for every method.
Setting $\mathcal{A}_t = \{e\}$ gives the single-agent reference $\hat{y}^{e,\mathrm{solo}}_t$.

\subsection{Task Definitions}
\label{sec:bench:tasks}

\paragraph{T1 --- Cooperative point cloud reconstruction.}
\emph{Input:} ego sweep $P^e_t$, messages from all collaborators, and relative
poses. \emph{Output:} a point set $\hat{Q}^e_t \subset \mathbb{R}^3$ in the ego
frame over a fixed evaluation volume $V$ (an $L \times W \times H$ box centred
on the ego), optionally with per-point confidence.
\emph{Ground truth:} the union of all agents' sweeps in the window, registered
into the ego frame with surveyed poses and voxel-deduplicated at resolution
$\rho$, restricted to $V$. \emph{Why collaborative:} the target is the geometry
the ego cannot see, so the single-agent reference is bounded away from the
ground truth by construction; the task measures how much of that gap a method
recovers.
\todo{Fix $V$, $\rho$, and the registration tolerance; state whether dynamic
objects are excluded from the reconstruction ground truth.}

\paragraph{T2 --- Cooperative 3D object detection.}
\emph{Input:} as T1. \emph{Output:} a set of oriented 3D boxes
$\{(x,y,z,l,w,h,\theta,c,s)\}$ in the ego frame, with class $c$ and score $s$.
\emph{Ground truth:} the union of all agents' annotations transformed into the
ego frame, filtered to objects observed by at least one agent within the
evaluation range --- the same ground truth is used for the collaborative and
single-agent runs, so the reference is scored on objects the ego alone cannot
see.
\todo{State the evaluation range, the minimum-points-per-box visibility rule,
and the class list.}

\paragraph{T3 --- Cross-modal localisation.}
\emph{Input:} one query observation from a single modality --- a camera image
$I^{i,k}_t$, a radar return $R^i_t$, or a CSI snapshot $H^i_t$ --- plus a
reference representation of the environment (the ego's LiDAR BEV, or the
prior map, depending on the sub-protocol). \emph{Output:} the pose of the
querying agent, $\hat{T}^i_t$, as 3-DoF $(x, y, \psi)$ in the shared frame for
the map-referenced sub-protocol and as relative pose
$\hat{T}^{i \to e}_t$ for the ego-referenced one.
\emph{Why collaborative:} relative pose is the input every fusion operator in
T1, T2, and T4 depends on; this task measures the quality of that input in
isolation, per modality, rather than assuming ground-truth extrinsics.
\todo{Decide whether the CSI sub-protocol is evaluated per snapshot or per
short window, and fix the window length.}

\paragraph{T4 --- Communication-efficient collaborative perception.}
\emph{Input:} as T2, plus a bit budget $B$ bits per frame per link.
\emph{Output:} the T2 output produced under $|m^{i \to e}_t| \leq B$, together
with the measured message size. A method is characterised not by a single
number but by its \textbf{operating curve} over $B$, from the late-fusion
regime (box lists, $10^2$--$10^3$ bits) to the early-fusion regime (raw points,
$10^7$ bits). \emph{Why separate from T2:} T2 lets each method spend whatever
it natively spends, which makes its ranking partly a bandwidth ranking; T4
puts the cohort on a common budget axis.
\todo{Fix the budget grid $B$ and the accounting rule (payload only vs.\ payload
plus pose/metadata header).}

\paragraph{T5 --- Agent fingerprinting.}
\emph{Input:} a sequence of PHY-layer observations from one transmitter ---
CSI magnitude/phase across subcarriers, RSSI, and per-packet metadata --- over
a window of $W$ packets. \emph{Output:} an identity label over the $N$ enrolled
agents, with an open-set reject option for transmitters not in the enrolment
set. \emph{Why in a perception benchmark:} fusion operators weight messages by
claimed identity and pose; fingerprinting establishes how well identity can be
established from the physical layer alone, independent of any claim the message
makes about itself.
\todo{Fix $W$, $N$, the enrolment/query split, and whether the open-set
negatives come from held-out devices or held-out environments.}

\subsection{Splits and Metrics}
\label{sec:bench:splits}

\paragraph{Protocols.}
Two protocols are defined for every task.
\emph{Within-environment (WE)}: train and test in the same environments, split
by disjoint time segments with a guard interval of $\Delta$ seconds between
segments so that no test frame shares a temporal neighbourhood with a training
frame. \emph{Leave-one-environment-out (LOEO)}: train on all environments but
one, test on the held-out one, repeated over environments; the reported number
is the mean over folds with per-fold numbers in the appendix. LOEO is the
headline protocol --- WE measures fit, LOEO measures transfer.
\todo{Enumerate the environments, the frame counts per split, the guard
interval $\Delta$, and the frame stride used at evaluation.}

\paragraph{Metrics.}
\begin{itemize}
  \item \textbf{T1:} Chamfer distance (symmetric, in metres) and F-score at
        thresholds $\{0.1, 0.2, 0.5\}$\,m, computed inside $V$; plus
        \emph{coverage of the ego-occluded region} --- F-score restricted to
        the volume the ego alone does not observe, which is the part
        collaboration exists to fill.
  \item \textbf{T2:} AP at IoU $0.5$ and $0.7$ (BEV and 3D), reported overall
        and split into ego-visible and ego-occluded zones.
  \item \textbf{T3:} median and 90th-percentile position error (m) and heading
        error (deg); recall at $\{0.5, 1, 2\}$\,m for the retrieval-style
        camera and CSI sub-protocols; ATE over trajectories where the
        sub-protocol is sequential.
  \item \textbf{T4:} accuracy-vs-bandwidth curves (AP at IoU $0.7$ against
        measured bits/frame, log axis) summarised by the area under the curve
        over the budget grid, and by the smallest budget at which a method
        stays within $1$ AP point of its unconstrained T2 score.
  \item \textbf{T5:} closed-set accuracy and macro-F1; open-set AUROC and
        TPR at $1\%$ FPR.
\end{itemize}
Every task additionally reports the \textbf{single-agent reference} and the
\textbf{collaboration gain} (collaborative minus reference) in the same units,
and \textbf{bits/frame} for the message the method actually sends.
\todo{State the number of seeds per cell and whether the reported number is
mean or median across seeds; give the dispersion measure used in all tables.}

\subsection{Baselines}
\label{sec:bench:baselines}

Baselines are grouped by the modalities they consume. Within each group,
methods are ordered by fusion stage --- \emph{early} (raw measurements),
\emph{intermediate} (learned features), \emph{late} (task outputs) --- because
stage determines both message size and where a collaborator's evidence enters
the computation. All detectors in the LiDAR-only cohort share one backbone so
that differences are attributable to the fusion operator rather than to the
detector.
\todo{Name the shared backbone and state, for each method, whether the
checkpoint is retrained on this dataset or ported; give training schedules in
the appendix.}

\subsubsection{Single-agent references}
Per-modality, no messages exchanged: a LiDAR detector, a camera BEV detector,
and a radar detector for T2; single-sweep accumulation for T1; per-modality
regression/retrieval for T3.
\todo{Cite and specify the reference detectors.}

\subsubsection{LiDAR-only collaborative methods}
\begin{itemize}
  \item \textbf{Early fusion} --- transmit raw points, detect once over the
        union \cite{cooper}. Upper bound on information, upper bound on
        bandwidth.
  \item \textbf{Late fusion} --- transmit boxes, merge by NMS. Lower bound on
        bandwidth.
  \item \textbf{Intermediate fusion} --- transmit BEV features:
        attention-based \cite{opv2v,cobevt,v2xvit}, maxout \cite{fcooper},
        graph message passing \cite{v2vnet,disconet},
        alignment-corrected \cite{coalign}, and
        selective/handshake \cite{where2comm}, which additionally exposes a
        bandwidth knob and is therefore also a T4 baseline.
\end{itemize}
\todo{Fix the final list; cite each method at its original publication and note
the implementation used.}

\subsubsection{Multi-modal collaborative methods}
\begin{itemize}
  \item \textbf{Camera-only cooperative} --- per-agent BEV lifting followed by
        cross-agent fusion \cite{cobevt,coca3d}. Tests whether collaboration
        recovers the depth that a single camera rig lacks.
  \item \textbf{LiDAR${+}$camera} --- per-agent sensor fusion into a shared BEV
        \cite{bevfusion} composed with each collaborative operator above.
  \item \textbf{Heterogeneous agents} --- collaboration between agents with
        \emph{different} sensor suites \cite{hmvit,heal}, which is the
        configuration the dataset is built for and which the LiDAR-only cohort
        cannot express.
  \item \textbf{Radar-inclusive} --- LiDAR${+}$4D-radar and radar-only
        cooperative detection \cite{v2xr}.
  \item \textbf{RF/CSI} --- for T3 and T5 only: classical fingerprinting
        ($k$-NN over CSI amplitude), AoA/ToF geometric estimation
        \cite{spotfi}, and a learned CSI encoder.
\end{itemize}
\todo{Fix the final list; state which multi-modal baselines are ported and
which are constructed here by composing a per-agent fusion frontend with an
existing collaborative operator --- constructed baselines must be labelled as
such in every table.}

\subsubsection{Baseline summary table}
\todo{Table: one row per baseline; columns = modalities consumed, fusion stage,
message content, nominal bits/frame, tasks it is evaluated on, source of the
implementation. Group the rows into the three blocks above with
\texttt{\textbackslash midrule} separators.}

\subsection{Main Results}
\label{sec:bench:results}

Results are reported in two tables that are never merged: Table~A for the
LiDAR-only cohort, Table~B for the multi-modal cohort. Each carries the
single-agent reference for its own input budget as its first row, so that
collaboration gain is read down a column rather than inferred across tables.
Both are reported under LOEO, with WE in the appendix.

\todo{Table A (LiDAR-only, per task): single-agent reference, then early /
intermediate / late fusion. Columns: T1 Chamfer $\downarrow$, T1 F@0.2
$\uparrow$, T1 occluded-F $\uparrow$, T2 AP@0.5 / AP@0.7 $\uparrow$, T2
occluded-zone AP $\uparrow$, bits/frame.}

\todo{Table B (multi-modal, same columns plus a modality column): camera-only,
LiDAR${+}$camera, heterogeneous, radar-inclusive; T3 in its own table with
per-modality rows (camera / radar / CSI) against the same targets.}

\todo{Headline numbers to state in prose: (i) collaboration gain over the
single-agent reference per cohort, and how much of it lands in the
ego-occluded zone; (ii) whether adding a modality or adding an agent buys more,
at matched bits/frame; (iii) the bandwidth at which the T4 curves cross.}

\subsection{Ablations}
\label{sec:bench:ablations}

\begin{enumerate}
  \item \textbf{Modality on/off} (multi-modal cohort only): each modality
        removed in turn, at fixed architecture and fixed number of agents.
  \item \textbf{Number of cooperating agents}: $|\mathcal{A}_t| = 1, 2, 3, \dots$,
        reported as a curve per cohort; establishes where the gain saturates.
  \item \textbf{Fusion stage at matched bandwidth}: early / intermediate / late
        compared at equal bits/frame rather than at their native sizes, which
        is the only comparison in which stage is the sole variable.
  \item \textbf{Pose source}: ground-truth extrinsics vs.\ poses estimated by
        the T3 baselines --- quantifies how much of the reported collaborative
        accuracy is contingent on perfect alignment.
  \item \textbf{Sensor placement and range}: performance vs.\ inter-agent
        distance and vs.\ relative viewpoint angle.
\end{enumerate}
\todo{Fix the agent-count grid and the matched-bandwidth operating point; state
which ablations run on WE and which on LOEO.}

\subsection{Qualitative Results}
\label{sec:bench:qualitative}

\todo{Figure 1: cooperative vs.\ single-agent reconstruction of the same
frame, side by side, with the ego-occluded region shaded in both --- the gap
in the single-agent panel is the claim.}

\todo{Figure 2: a detection frame where collaboration recovers a fully occluded
object, with each collaborator's contribution to the fused BEV shown
separately.}

\todo{Figure 3: LiDAR-only vs.\ multi-modal on the same frame, chosen so the
extra modality is what makes the difference (e.g.\ a low-reflectivity object
that the camera or radar branch recovers).}

\todo{Figure 4: a failure case --- a frame where collaboration does not help or
actively hurts (large relative pose error, poor viewpoint overlap, or a
disagreeing collaborator) --- and a one-sentence diagnosis.}
